\section{Related record-reduction methods}
\label{sp:app:reduction}

This appendix reviews the reduction families summarized in \autoref{sp:sec:reduction_methods},
in the order an analyst might try them.

\subsection{Random sampling}

The simplest reduction draws records at random. Under simple random sampling without replacement each
record has inclusion probability $n/N$ and design weight $N/n$, and the Horvitz--Thompson estimator
recovers population totals without bias \citep{horvitz1952, cochran1977, sarndal1992}. A random
subsample preserves the conditional distribution of the data in expectation, but it does not
reproduce administrative totals, so it has to be calibrated after the draw. Random selection is also
the reference that the data-reduction literature treats as the genuine test: informativeness-based
schemes are measured against it \citep{ma2015leveraging}, and pruning methods in deep learning often
fail to beat it at high reduction ratios \citep{guo2022deepcore}. In our setting it is the
reduce-first, calibrate-after baseline, a random draw followed by gradient-descent reweighting to the
targets. Replacing the gradient reweighting step with raking gives the classical sampling-then-raking
workflow; it is directly comparable on a raking-compatible categorical margin subset, but not on the
full mixed count-and-dollar target surface without changing the problem into generalized calibration.

\subsection{Raking, survey-weight sampling, and integerisation}

A second approach calibrates first and reduces second. Given fractional weights from a calibration
step, survey-weight sampling selects records with probability proportional to those weights, which
turns a reweighted file into a finite, integer-weighted population that approximately preserves both
the calibrated totals and the distribution. This is unequal-probability,
probability-proportional-to-size sampling in the survey tradition \citep{hansen1943, horvitz1952},
applied to microsimulation as integerisation. The upstream calibration can be IPF/raking, GREG, or a
numerical optimizer; the downstream question is how to convert fractional weights into a finite file.
\citet{lovelace2013trs} catalogue the common integerisation methods and show that
truncate-replicate-sample reproduces the targets most accurately while fixing the reduced population
to an exact size; earlier top-up schemes are due to \citet{ballas2005}. Because the selection
probabilities are the calibration weights, survey-weight sampling presupposes a calibration step and is
the calibrate-then-reduce member of the literature most adjacent to our setting. It is informed
selection, but it prioritises records by the population mass each carries rather than by their
contribution to target fit: a record is retained with probability proportional to its prior weight, so
the step keeps records that represent more population, not records that help match a target the fit would
otherwise miss. The administrative targets enter only through the upstream weights, and the sampling
preserves that calibrated distribution instead of re-optimising against the targets. We therefore treat
it as a related literature method rather than a
target-fitting baseline run here, and note it as a natural calibrate-then-reduce comparator for a future
sweep (\autoref{sp:sec:future_work}). A raking-then-TRS variant is the corresponding classical baseline
on a categorical margin surface, but it does not naturally extend to the full mixed production target
surface without replacing IPF with a broader calibration solver.

\subsection{Balanced sampling}

Balanced sampling selects a fixed-size sample whose Horvitz--Thompson estimates reproduce auxiliary
totals exactly or nearly exactly. The cube method of \citet{deville2004cube} is the standard
construction: it chooses inclusion indicators while preserving prescribed balancing equations, then
uses the resulting sample weights for estimation. This makes balanced sampling a target-aware
selection method rather than a post-sampling calibration method, and therefore a relevant published
neighbour of $L_0$ selection. Its fit to the present problem is not exact. Classical balanced sampling
usually treats the auxiliary totals as design constraints with predetermined inclusion probabilities,
whereas the production calibration surface here mixes categorical counts, dollar totals, and
hierarchical relationships and then fits positive weights after selection. Still, it is an important
reference point for the idea that the sample itself, not only its weights, can be chosen to preserve
known totals.

\subsection{Convex sparse calibration}

A final family keeps the continuous optimization setup but replaces the direct retained-record count
with a convex sparsity surrogate. An $L_1$ penalty on fitted weights, implemented with proximal
soft-thresholding, can drive weights exactly to zero while preserving the same calibration loss, and it
is tractable on the full target surface. Its limitation is structural rather than computational: the
single penalty couples how many records survive with how much the surviving weights shrink, so its
tuning parameter is not a direct retained-record-count control in the way the expected $L_0$ norm is. We
therefore read it as a penalty-based analogue of the $L_0$ selector that this study does not implement,
and one we expect to make a weaker sampler for that reason; Section~\ref{sp:sec:l1_background} develops the
mechanism and the coupling argument.

\subsection{Informed selection}

These methods share the idea that a carefully chosen weighted subset can stand in for a larger
dataset, but they reach it differently: random sampling ignores the targets and repairs the gap
afterwards, survey-weight sampling inherits the targets from a prior calibration, balanced sampling
tries to preserve auxiliary totals through the sample design, and convex sparse calibration remains in
the same optimizer but changes the penalty. A fourth tradition makes selection itself the calibration:
combinatorial optimisation chooses a combination of records, with integer multiplicities, to match the
benchmark totals directly, searched with metaheuristics such as simulated annealing
\citep{kirkpatrick1983, williamson1998, voas2000}. It is the closest precedent for target-informed
selection, but it is discrete and gradient-free; we note it as related work rather than include it in
the matched-budget comparison. The method studied here also pursues the targets at selection, but with
a continuous, gradient-based relaxation that fits the selection and the weights jointly. The next
subsection sets out that relaxation.

